\chapter{Introduction}

% Accroche & introduction du sujet

\begin{quote}
    \textit{Mais dès l'instant où le critère de l'authenticité s'avère inapplicable à la production artistique, la fonction sociale entière de l'art s'en trouve bouleversée. Au lieu de se fonder sur le rituel, elle se fonde sur une autre praxis : en l'occurrence, sur la politique.}
    
    Walter Benjamin, \textit{L'\oe uvre d'art à l'époque de sa reproductibilité technique}, 1939.
\end{quote}

L'acte de copie d'un objet culturel est la première étape à son partage, sa conservation et sa réinterprétation. Si une langue a besoin de locuteurs pour pouvoir évoluer, alors un objet culturel qui n'est pas copié ne peut pas être nommé comme un \guillemetleft objet culturel \guillemetright ~, tout comme le latin ne peut pas être considéré comme une \guillemetleft langue vivante \guillemetright. Un premier exemple de l'importance de la copie et par extension du piratage dans notre patrimoine est \textit{Le Voyage dans la Lune} réalisé en 1902 par Georges Méliès. La plupart des copies officielles de ce film ont disparu, ou ont été détruites par Méliès lui-même alors qu'il ne pouvait plus les entreposer à cause de la fermeture de son studio. Les chercheurs et restaurateurs ont pu reconstituer et préserver son œuvre grâce à de très nombreuses copies qui étaient soit piratées, soit confisquées \footnote{Marie-France Briselance et Jean-Claude Morin, \textit{Grammaire du cinéma}, Paris, Nouveau Monde, coll. « Cinéma », 2010, 588 p. (ISBN 978-2-84736-458-3), p. 111.}.

Un deuxième exemple, peut-être le plus connu, est celui du film \textit{Nosferatu le vampire} réalisé en 1922 par Friedrich Wilhelm Murnau. Florence Stoker, veuve de Bram Stoker, l'auteur du roman \textit{Dracula}, intente à Murnau un procès pour plagiat que ce dernier perd. Toutes les copies ainsi que les négatifs sont détruits par deux fois en 1925 et 1928. Rapidement, d'anciennes copies cachées refont surface - et le film survit grâce à des projections confidentielles jusqu'en 1981 où plusieurs copies partielles permettent de le restaurer au plus près de l'archétype du travail de Murnau \footnote{Jean Pierre Piton, \textit{L'éternel retour de Dracula}, L'Écran Fantastique no 130 janvier 1993 p. 21.}.

La numérisation de nos outils professionnels, de nos sources de divertissement et de nos modes de communication, généralise la capacité de copie des objets culturels qui peuplent nos vies. Cette propriété préexistait déjà, avec, par exemple, la K7 VHS à bande magnétique, qui permet d'enregistrer un signal analogique. Cependant, dans un objet numérique, la propriété de copie est parfaitement intégrée au fonctionnement du fichier et de son système d'accès par défaut, l'ordinateur, produisant un dilemme numérique entre la liberté du partage d'une information réplicable gratuitement et la propriété intellectuelle littéraire et artistique \footnote{\cite{national_research_council_us_digital_2010}}.

La nouvelle facilité de copie et d'échange de l'information grâce à la généralisation de l'informatique bon marché et de l'accès au web domestique a rapproché des individus qui se sont organisés en communautés culturelles,  partageant des références culturelles et les fichiers qui les documentent, devenant ainsi des communautés patrimoniales.

Les communautés patrimoniales \footnote{~Convention de Faro sur la valeur du patrimoine culturel pour la société, 2005} qui nous intéressent dans ce travail sont celles qui décident d'embrasser la copie et le partage de l'information dans sa forme la plus radicale, en privilégiant les \guillemetleft communs numériques \guillemetright. Elles peuvent choisir de produire et d'entretenir leur propre contenu à la manière d'un Wikipédia, devenant des \textit{communautés libres}, ou elles peuvent choisir de faire fi des règles et des lois en devenant des \textit{communautés patrimoniales pirates}\footnote{\cite{bermes_lecran_2024}}. Si les communautés libres essaient de se différencier le plus possible des communautés pirates, les individus les composant évoluent entre les deux, puisant dans les ressources des communautés libres pour ensuite compléter leurs besoins avec l'aide du piratage.

Nous étudierons ici les liens, les émotions, et les mouvements de va-et-vient entre les individus qui composent les communautés pirates, les objets numériques piratés, et les lieux virtuels où se construisent ces communautés. Notre premier but est d'étendre le concept de communautés patrimoniales aux pirates culturels, ensuite de comprendre l’appréhension du public du \guillemetleft dilemme numérique de la propriété intellectuelle \guillemetright, et enfin d'étudier les pratiques émergentes de partage et de conservation numériques qui mobilisent ces communautés patrimoniales pirates.

\section*{Présentation bibliographique}

Les travaux sur le piratage se répartissent sur deux axes bibliographiques principaux. Le premier se concentre sur les industries de la propriété intellectuelle \footnote{\cite{liang_piratage_2014}}, essentiellement dans le domaine du marché pharmaceutique international. 

Le second axe est celui du travail quantitatif sur le piratage de biens culturels.Le texte se focalise sur le problème des pertes potentielles pour le secteur de l'édition, du jeu vidéo ou pour les canaux de diffusion sportif. Le piratage étant un sujet actuel, une grande partie des sources à disposition sont composées d'articles de presse toujours rédigés en deux parties, la première évalue la part du manque à gagner du piratage pour ensuite évaluer l'efficacité des mesures répressives mises en place au niveau national. Nous essaierons de nous détacher de ces deux approches suffisamment cartographiées pour comprendre les \guillemetleft carrières de pirate \guillemetright ~des personnes composant ces communautés patrimoniales \footnote{\cite{galluzzo_longevite_2021}}.

\subsection*{Le \guillemetleft piratage \guillemetright~: entre enjeux scientifiques et communication  }%je ne suis pas satisfait de ce titre

Dans ce travail, nous utiliserons le champ lexical de la piraterie historique par souci de simplicité et d’accessibilité, cependant il me semble important d'insister sur le fait que produire une copie d'un objet numérique n'est pas assimilable à du \textit{piratage}. La compréhension populaire du partage non autorisé de contenu culturel est influencée par une communication industrielle relayée dans les milieux médiatique et institutionnel. Cette approche associe le \textit{partage non autorisé} au mot \textit{piratage}, en comparant le partage d'informations à la contrefaçon, au vol et au piratage informatique. Ce discours est au cœur de stratégies de communication de plusieurs institutions publiques, dont l'exemple le plus connu est le clip \textit{Voler une voiture ? Jamais !} publié en 2004 par la FACT (Federation Against Copyright Theft) et la Motion Picture Association of America, comparant directement le partage de contenu culturel sur Internet à plusieurs actes de vol avec effraction. 

Le piratage est un sujet médiatique récurrent et propice à la communication. En plus des rapports des institutions étatiques, l'exemple le plus visible est celui de Netflix, qui a fortement communiqué sur son modèle économique décrit comme une \guillemetleft alternative au piratage \guillemetright \footnote{Sophia Harris, \textit{Netflix's anti-piracy team aims to make stealing content uncool}, CBC News, 2017}. Un autre exemple est la communication de Gabe Newell, cofondateur et président de l'entreprise Valve - absolument dominant dans le domaine de la distribution du jeu vidéo dématérialisé sur PC.

\begin{quote}
    \textit{Nous pensons qu'il y a un malentendu fondamental sur le piratage. Le piratage est presque toujours un problème de service et non un problème de prix », a-t-il déclaré. « Si un pirate propose un produit partout dans le monde, 24 heures sur 24, 7 jours sur 7, que vous pouvez acheter depuis votre ordinateur personnel, et que le fournisseur légal dit que le produit est verrouillé par région, qu'il sera disponible dans votre pays 3 mois après la sortie aux États-Unis et qu'il ne peut être acheté que dans un magasin de briques et de mortier, alors le service du pirate a plus de valeur \footnote{Entrevue avec Gabe Newell dans \textit{The Cambridge Student}, via le billet de blog sur le site \textit{The Escapist} par Greg Tito le 28/11/2011} ».}
\end{quote}

La pratique du téléchargement illégal obéit à des stratégies et à des questions plus complexes que la seule question de l'accessibilité en ligne d'une ressource convoitée - et ne peut être résumée à une manne de consommateurs attendant un acteur entrepreneurial prêt à conquérir ce nouveau marché. 

Le mot \guillemetleft  Piratage \guillemetright~, ainsi que l'iconographie de la piraterie historique, est également un point de rendez-vous, ou plus rarement un repoussoir - pour les communautés qui investissent le web comme un espace d'échange de biens culturels. L'exemple le plus partagé est le nom et l'iconographie du site \textit{The Pirate Bay}, dont l'icône est un trois-mâts noir, dont la voile centrale montre une variation du \textit{Jolly Roger} remplaçant le crâne par une cassette en inscrivant en dessous dans une police gothique le nom du site de partage de torrents le plus connu des internets.  L'utilisation d'autres termes comme \textit{partage illégal du contenu culturel, reproduction interdite ou non autorisée, partage d'informations avec son voisin} est une alternative qui offre une représentation plus mesurée et plus proche de la réalité des faits \footnote{Richard Stalman,\textit{Quelques expressions trompeuses qu'il vaut mieux éviter d'employer}, dans \guillemetleft Libre enfants du savoir numérique \guillemetright, p.359-364 }. Au-delà de cette analyse légaliste et quantitative, la sphère académique des sciences humaines et sociales s'est intéressée à la popularité de plateformes telles que Napster \footnote{Moulier-Boutang, \textit{Droits de propriété intellectuelle, terra nullius et capitalisme cognitif}. Multitudes,n°41(2), p66-72, 2010} ou The Pirate Bay . Ces recherches universitaires cherchent à comprendre la psychologie des individus en examinant les discours et valeurs associés au piratage. La méthodologie de notre travail explore et combine les approches sociologiques et les études archivistiques des pratiques de conservation et de partage des objets numériques. 

Le sujet qui reste le plus étudié est la place du piratage dans les économies des pays en développement, à la fois ateliers du monde du capitalisme planétaire, mais aussi envisagé comme une manne de consommateurs potentiels des divertissements du nord économique. Des revues de science de l'information, de communication, de management et d'économie étudient l'organisation et les chemins de la distribution du contenu culturel des pays du Sud économique comme par exemple à Malégaon en Inde \footnote{Lawrence Liang,\textit{ Piratage, créativité et infrastructure : repenser l’accès à la culture}, Tracés. Revue de Sciences humaines, 2014}, ou au Cameroun \footnote{Mattelart, T.\textit{ Chapitre 1. Piratages audiovisuels et réseaux de la mondialisation par le bas}. Dans \og Piratages audiovisuels : Les voies souterraines de la mondialisation culturelle \fg, p. 27-52, éditions De Boeck Supérieur, 2011}.

\subsection*{Définir la communauté patrimoniale pirate}

L'objet numérique et son partage mobilisent des concepts et des catégories juridiques diverses qu'il est important de définir. Cet objet numérique s'incorpore dans le concept de \textit{patrimoine documentaire immatériel}, promu et adopté par la Convention de l'UNESCO pour que ses États membres puissent intégrer au droit des définitions et des outils politiques de promotion de la diversité culturelle, y compris de la culture numérique et de la culture du web. Si un patrimoine est immatériel, c'est alors qu'il est constitué d'objets qui ont un sens particulier pour un groupe de personnes, le patrimoine immatériel étant donc la relation entre ces \textit{objets} et les groupes de personnes qui leur donnent un sens et une fonction. Quand un objet est la cause principale de lien dans un groupe (et non plus une conséquence de ce lien), nous pouvons alors parler de l'apparition d'une \textit{communauté patrimoniale}. Ce terme de \textit{communauté patrimoniale} entre dans le vocabulaire institutionnel en 2005 avec la signature de la convention européenne de Faro sur la valeur du patrimoine culturel pour la société. Cette convention inscrit l'action européenne dans le domaine du patrimoine avec un nouveau paradigme, considérant autant les objets et les lieux que les usages, les personnes et les interactions qui en résultent.

L'émotion négative est une composante puissante dans la définition de ce qu'est le patrimoine (peur de voir disparaître, de voir changer, de voir trahir un objet), les communautés se soudent avec des émotions comme la colère, la sédition, la rébellion face à l'injustice \footnote{\cite{fabre_emotions_2013}}, en invoquant des référentiels politiques ou institutionnels.  Les communautés patrimoniales peuvent entrer dans des dynamiques d'actions de sauvegarde de leur patrimoine, parfois avec la volonté d'aider les acteurs institutionnels ou entrepreneuriaux (une restauration d'un bâtiment par exemple) - ou en entrant dans un rapport de force avec ces mêmes acteurs (la mobilisation des mouvements écologistes face à des travaux d'infrastructures comme les bassines de Sainte-Soline en est un bon exemple). 

Si le partage non autorisé de biens culturels n'a rien de nouveau, le tournant numérique de nos sociétés a engendré deux conséquences. La première est la massification des échanges, des capacités de stockage et de copie des individus - la deuxième est la capacité nouvelle des individus à s'affranchir du temps et de l'espace pour entrer en relation, donnant ainsi naissance à des communautés patrimoniales plus nombreuses et beaucoup plus développées dès la démocratisation de l'accès au web. Quand une communauté patrimoniale se forme autour d'objets ou de pratiques \textit{illégales} nous pouvons alors utiliser le terme de communautés patrimoniales \textit{pirates}. Notre travail a pour objectif de mieux comprendre les relations entre l'objet, l'individu et la communauté dans ce cadre de sortie des référentiels institutionnels et entrepreneuriaux - le but des pirates n'étant pas de trouver une légitimité aux yeux de l'État ou des entreprises, mais bien de s'en affranchir.  

\subsection*{De l'individu à la communauté, progresser de la carrière pirate aux dépôts communautaires en ligne}

Étudier des individus qui piratent implique de se plonger dans leurs histoires personnelles, et les origines de leur montée en compétence progressive - passant de néophyte à amateur, puis pour certains à des professionnels du piratage, ce qu'Anthony Galluzzo, professeur en science de gestion à l'université de Saint-Étienne, nomme une \textit{carrière pirate}\footnote{\cite{galluzzo_longevite_2021}}. Cette acquisition de compétences techniques est parallèle à une éducation, un apprentissage social et politique - c'est en combinant les deux que les pirates génèrent une sociabilité particulière s'appuyant autant sur des éléments techniques (comment accéder à une ressource, qui possède un tutoriel pour tel outil logiciel, qui possède une copie d'un objet numérique rare que je cherche) que sur des émotions patrimoniales (ce site va disparaître - comment aider ma communauté à le préserver). 

Ce lien a été l'objet d'étude de la littérature anglo-saxonne des \textit{fan studies} \footnote{\cite{taylor_save_nodate}} \footnote{:\cite{de_kosnik_rogue_2016}}. Ces études, peu nombreuses mais présentes dans l'actualité scientifique, se concentrent sur les cultures du web (des fanfictions, des œuvres d'art) et la manière dont les personnes qui vivent ces cultures comprennent et organisent la mémoire de ce qui a été produit au sein de leur communauté. Notre travail est à la fois proche car nous y étudions une culture \textit{nativement numérique}, et en même temps certaines différences fondamentales de la situation des communautés patrimoniales pirates ne permettent pas de simplement plaquer les analyses des \textit{fan studies} sur ce sujet. 

L'un des paramètres ayant le plus d'influence sur l'organisation des communautés patrimoniales pirates et leurs membres est \textit{l'organisation anarchique de guérilla}. Comme pour Wikipédia, la plupart des dépôts communautaires pirates n'obéissent pas à une structure verticale et rigide. Un noyau dur de membres reconnus pour leur implication et leur ancienneté dans le projet permet de maintenir l'infrastructure de base du dépôt communautaire - ce à quoi s'ajoutent des membres moins investis qui offrent leur aide et leurs conseils de manière ponctuelle. Un troisième cercle est la communauté large autour du dépôt, des usagers réguliers offrant un soutien financier, qui propagent la bonne adresse à l'aide du bouche à oreille sur les forums, ou simplement qui témoignent de l'importance du dépôt pirate dans leur vie de tous les jours. 

Le terme de guérilla fait référence à la manière dont ces dépôts, leurs membres et leurs usagers s'organisent pour mettre en place un tiers-lieu à la fois efficace dans sa mission (transmettre et archiver des données) sans se faire repérer par les éditeurs ou les autorités de protection des droits d'auteur \footnote{Martin Paule Eve, Lessons from the Library: Extreme Minimalist Scaling at Pirate Ebook Platforms, dans \textit{Digital Humanities Quarterly} Vol 16 n°2, 2022}. Pour cela, les dépôts pirates doivent à la fois avoir une visibilité de surface minimale (ne pas être vus), sans que le trafic internet puisse être retracé à l'hébergeur du dépôt (ne pas être visé), avec une résilience aux attaques et aux fermetures d'hébergeur (ne pas être touchés).

\section*{Méthodologie et problématique}
%pas encore utilisé
Notre méthodologie s'inspire à la fois des \textit{fanstudies}, de la sociologie de la déviance et des études en archivistique et en sciences de l'information et des bibliothèques. Notre but est de rapprocher les expériences personnelles des personnes qui piratent et l'étude de l'architecture organisationnelle des bibliothèques pirates pour comprendre les éléments et les formes d'organisation dans les communautés patrimoniales pirates. \newline

Une partie de mon étude se base sur l'étude des archives et des forums des communautés patrimoniales pirates. Une partie a fait l'objet d'un archivage, parfois d'une patrimonialisation et d'une curation - c'est par exemple le cas des archives du web de la BnF ou des collections de The Internet Archive. Une autre partie de mes sources a été directement archivée en utilisant l'outil d'instantané (snapshot) Zotero.
\begin{itemize}
    \item Mes premières sources sont les archives du web de la BnF, en particulier les archives du site Skyblog. Grâce à Marina Hervieu, chargé de projet de recherche à la BnF, j'ai pu de manière fiable explorer les skyblogs ayant dans leurs pseudonymes des mots laissant penser à une pratique régulière du piratage \footnote{pirate, piratage, téléchargement, télécharger, torrenting, torrent, fansub, mininova}. En utilisant cette méthode, j'ai très rapidement atteint une saturation sémantique, où les exemples se ressemblent beaucoup sans vraiment permettre de creuser les situations que j'avais étudiés.
    
    \item Une partie de mes sources est de monitorer les archives et l'actualité de forums de communautés patrimoniales pirates, particulièrement sur Reddit, Twitter/X et Discord. Les deux sites les plus intéressants à étudier ont été discord et Reddit car Twitter/X est une plateforme qui ne propose de former des "communautés" que depuis très récemment. 
    
    \subitem Discord est un système de messagerie communautaire où j'ai suivi la communauté qui se nomme \textit{The Eye}. C'est un site communautaire, semblable à une association ayant pour but de préserver et de partager des pans de la culture numérique. 

    \begin{quote}
        \textit{Living without the knowledge of our past history, origin and culture is like a tree without roots. The Internet is a worldwide platform for sharing information. It is a community of common interests.
        Our mission at The Eye is to preserve pieces of digital history. We are digital librarians.}\footnote{Message de présentation du serveur discord de \textit{The Eye}}
    \end{quote}

    \subitem Reddit est également un site très instructif pour étudier les communautés patrimoniales pirates. Il est composé de \textit{subreddits}, des forums communautaires sur des sujets spécifiques. Il existe plusieurs subreddit sur les sujets du piratage de bien culturel et de la sauvegarde du patrimoine numérique. Si tous un subreddit ne peuvent pas être considéré comme une communauté patrimoniale, certaines communautés investissent des subreddits particuliers - c'est le cas de \textit{r/datahoarder} (des individus collectant de grands volumes de données numériques) qui est un carrefour de discussion et de mobilisation d'amateurs d'archives numériques et de patrimoine dématérialisé. 

    \subitem Twitter/X est un endroit un peu plus difficile d'approche pour notre travail, il s'agit d'une application de microblog, où les gens peuvent se suivre de manière individuel sans avoir l'option d’agréger les utilisateurs en des communautés\footnote{Il existe depuis la reprise de l'application par Elon Musk un onglet communauté, mais la perte de popularité de Twitter couplet à l'ajout de cette fonction relativement peu utilisée fait que Twitter reste une application de microblog individualisée. }. Les sources que j’agrège de Twitter ne sont pas représentatives d'une communauté patrimoniale, cependant certains sujets transversaux peuvent mobiliser un grand groupe d'utilisateur - créant ainsi un effet de mode et d’entraînement permettant d'étudier l'impacte de certains évènements sur la création de communautés patrimoniales pirates. L'exemple le plus récent dont je peux témoigner est la mort du réalisateur américain David Lynch. Les \textit{twittos}\footnote{Twittos : surnom vulgaire donné à un utilisateur intensif du réseau Twitter/X} et fans endeuillé ont commencé à mettre à partager illégalement les films et séries et documentaires du réalisateur. 
    
    \end{itemize}

Dans le temps qui m'a été imparti, j'ai réalisé cinq entretiens avec une sélection de personnes. Mon but était de pouvoir fournir une diversité en termes d'âge, de genres, de pratiques de piratage. Ces entretiens ont été pour la plupart réalisés en tête-à-tête avec la personne, sauf pour les personnes trop lointaines et pour des personnes comme Thomas ou Élise où un entretien en visioconférence a été préféré. Une captation audio a été systématiquement réalisée afin de produire une transcription en utilisant le modèle Whisper d'OpenAI.\newline

Tableau des profils des différents entretiens.\\
\begin{tabular}{|p{2cm}|p{2cm}|p{2cm}|p{2cm}|p{2cm}|}

\hline
   Prénom  & Âge & Profession & Méthode & Résumé\\
\hline
    Adrien & 23 ans & Étudiant en informatique & Torrenting, Téléchargement & Diplôme en informatique \\
\hline
    Théo & 27 ans & Ingénieur informaticien & Torrenting, Serveur & Diplôme en informatique\\
\hline
    Hilaire & 26 ans & Étudiant en sociologie & Téléchargement, forums & Diplôme en communication et en sociologie \\
\hline 
    Zakary & 25 ans & Étudiant en communication & Téléchargement, communauté & Diplôme en histoire et communication \\
\hline
    Elise & 23 ans & Étudiante en communication & Téléchargement, communauté & licence langue et édition, master de communication \\
\hline
    Siham & -- ans & -- & - & - \\
\hline
\end{tabular}\\

J'ai tenté de représenter équitablement les genres, mais un temps insuffisant entraîne un biais persistant lié à l'âge. La plupart des répondants sont des personnes de ma génération (nées entre 1995 et 2003). J'ai cependant réussi à mêler des profils de personnes proches de l'informatique (la majorité du corps social qui m'entoure) et d'autres qui piratent sans pouvoir se reposer sur des connaissances techniques aussi poussées. 
Malgré une cohorte limitée, la saturation sémantique des réponses de mes enquêtés a été relativement rapide. La plupart de mes répondants ont une pratique du piratage privilégiant l'accès le plus simple et le plus rapide aux ressources, homogénéisant rapidement les réponses aux questions. Les entretiens individuels sont difficilement des portes d'entrée vers les communautés patrimoniales pirates - qui sont par définition des œuvres collectives ; ils restent cependant très intéressants pour comprendre la manière dont des individus interagissent avec ces communautés sans pour autant en faire partie.

